\ifdefined\gkver\else\def\gkver{student}\fi
\documentclass[\gkver]{gaokaozhenti}
\begin{document}

\chapter{1992年湖南云南海南}

\begin{enumerate}
\item
%% number: 1
%% paperName: 1992年湖南云南海南高考第 1 题 3 分
%% typeId: 1
%% type: 单选题
%% chapter: 磁场
%% point: 左手定则
%% method: 无
%% score: 3
%% degree: 650
%% duplicateId: 0
%% body:
两根互相平行的长直导线位于图中纸面内，导线中通有大小相等、方向相反的电流，如图所示。导线 a、b 所受的安倍力 $F_{a}$、$F_{b}$ 的方向是\xzanswer{A}
%% MISSING-IMAGE: 题干图（原 JSON 无 <img>），缺图待补；补图后用 \onepicture 插入
\fourchoices[answer=A]
{$F_{a}$ 向左、$F_{b}$ 向右}
{$F_{a}$ 向右、$F_{b}$ 向左}
{两力都垂直纸面，$F_{a}$ 向里、$F_{b}$ 向外}
{两力都垂直纸面，$F_{a}$ 向外、$F_{b}$ 向里}
%% answer: A
%% memo:
\memoanswer{
本题原卷及材料中未提供详解，待补充。
}

\item
%% number: 2
%% paperName: 1992年湖南云南海南高考第 2 题 3 分
%% typeId: 1
%% type: 单选题
%% chapter: 运动和力
%% point: 超重失重
%% method: 无
%% score: 3
%% degree: 850
%% duplicateId: 0
%% body:
用一根细绳将一重物吊在电梯的天花板上。在下列四种情况中，绳的拉力最大的是\xzanswer{C}
\fourchoices[answer=C]
{电梯匀速上升}
{电梯匀速下降}
{电梯加速上升}
{电梯加速下降}
%% answer: C
%% memo:
\memoanswer{
本题原卷及材料中未提供详解，待补充。
}

\item
%% number: 3
%% paperName: 1992年湖南云南海南高考第 3 题 3 分
%% typeId: 1
%% type: 单选题
%% chapter: 电场
%% point: 库仑定律
%% method: 无
%% score: 3
%% degree: 750
%% duplicateId: 0
%% body:
两个放在绝缘架上的相同金属球，相距 $d$，球的半径比 $d$ 小得多，分别带有 $q$ 和 $3q$ 的电荷，相互斥力为 $3F$。现将这两个金属球接触，然后分开，仍放回原处，则它们的相互斥力将变为\xzanswer{D}
\fourchoices[answer=D]
{0}
{$F$}
{$3F$}
{$4F$}
%% answer: D
%% memo:
\memoanswer{
本题原卷及材料中未提供详解，待补充。
}

\item
%% number: 4
%% paperName: 1992年湖南云南海南高考第 4 题 3 分
%% typeId: 1
%% type: 单选题
%% chapter: 原子和原子核
%% point: 核聚变
%% method: 无
%% score: 3
%% degree: 850
%% duplicateId: 0
%% body:
太阳辐射能量主要来自太阳内部的\xzanswer{D}
\fourchoices[answer=D]
{化学反应}
{放射性衰变}
{裂变反应}
{热核反应}
%% answer: D
%% memo:
\memoanswer{
本题原卷及材料中未提供详解，待补充。
}

\item
%% number: 5
%% paperName: 1992年湖南云南海南高考第 5 题 3 分
%% typeId: 1
%% type: 单选题
%% chapter: 气体性质
%% point: 分子动理论
%% method: 无
%% score: 3
%% degree: 850
%% duplicateId: 0
%% body:
花粉在水中做布朗运动的现象说明\xzanswer{B}
\fourchoices[answer=B]
{花粉的分子在做激烈的热运动}
{水分子在做激烈的热运动}
{水分子之间是有空隙的}
{水分子之间有分子力作用}
%% answer: B
%% memo:
\memoanswer{
本题原卷及材料中未提供详解，待补充。
}

\item
%% number: 6
%% paperName: 1992年湖南云南海南高考第 6 题 3 分
%% typeId: 1
%% type: 单选题
%% chapter: 运动和力
%% point: 牛顿定律的应用
%% method: 无
%% score: 3
%% degree: 550
%% duplicateId: 0
%% body:
从地面竖直上抛一小球。设小球上升到最高点所用的时间为 $t_{1}$，下落到地面所用的时间为 $t_{2}$。若考虑到空气阻力的作用，则\xzanswer{B}
\fourchoices[answer=B]
{$t_{1} > t_{2}$}
{$t_{1} < t_{2}$}
{$t_{1} = t_{2}$}
{因不知速度和空气阻力的关系，故无法断定 $t_{1}$、$t_{2}$ 哪个较大}
%% answer: B
%% memo:
\memoanswer{
本题原卷及材料中未提供详解，待补充。
}

\item
%% number: 7
%% paperName: 1992年湖南云南海南高考第 7 题 3 分
%% typeId: 1
%% type: 单选题
%% chapter: 电磁感应
%% point: 楞次定律
%% method: 无
%% score: 3
%% degree: 550
%% duplicateId: 0
%% body:
一个 N 极朝下的条形磁铁竖直下落，恰能穿过水平放置的固定小方形导线框，则\xzanswer{A}
%% MISSING-IMAGE: 题干图（原 JSON 无 <img>，磁铁经位置 1、2 图缺失），缺图待补；补图后用 \onepicture 插入
\fourchoices[answer=A]
{磁铁经过图中位置 1 时，线框中感应电流沿 abcd 方向，经过位置 2 时，沿 adcb 方向}
{磁铁经过 1 时，感应电流沿 adcb 方向，经过 2 时沿 abcd 方向}
{磁铁经过 1 和 2 时，感应电流都沿 abcd 方向}
{磁铁经过 1 和 2 时，感应电流都沿 adcb 方向}
%% answer: A
%% memo:
\memoanswer{
本题原卷及材料中未提供详解，待补充。
}

\item
%% number: 8
%% paperName: 1992年湖南云南海南高考第 8 题 3 分
%% typeId: 1
%% type: 单选题
%% chapter: 光学
%% point: 光的干涉
%% method: 无
%% score: 3
%% degree: 750
%% duplicateId: 0
%% body:
从点光源 L 发出的白光，经过透镜后成一平行光束，垂直照射到挡光板 P 上，板上开有两条靠得很近的平行狭缝 S$_{1}$、S$_{2}$，如图所示。在屏 Q 上可看到干涉条纹。图中 O 点是屏 Q 上与两狭缝等距离的一点，则\xzanswer{C}
%% MISSING-IMAGE: 题干图（原 JSON 无 <img>），缺图待补；补图后用 \onepicture 插入
\fourchoices[answer=C]
{干涉条纹是黑白的，O 是亮点}
{干涉条纹是黑白的，O 是暗点}
{干涉条纹是彩色的，O 是亮点}
{干涉条纹是彩色的，O 是暗点}
%% answer: C
%% memo:
\memoanswer{
本题原卷及材料中未提供详解，待补充。
}

\item
%% number: 9
%% paperName: 1992年湖南云南海南高考第 9 题 3 分
%% typeId: 1
%% type: 单选题
%% chapter: 磁场
%% point: 带电粒子在磁场中的运动
%% method: 无
%% score: 3
%% degree: 650
%% duplicateId: 0
%% body:
三个质子 1、2 和 3 分别以大小相等、方向如图所示的初速度 $v_{1}$、$v_{2}$ 和 $v_{3}$，经过平板 MN 上的小孔 O 射入匀强磁场，磁场方向垂直纸面向里。整个装置放在真空中，且不计重力。这三个质子打到平板 MN 上的位置到小孔 O 的距离分别是 $s_{1}$、$s_{2}$ 和 $s_{3}$，则\xzanswer{D}
%% MISSING-IMAGE: 题干图（原 JSON 无 <img>），缺图待补；补图后用 \onepicture 插入
\fourchoices[answer=D]
{$s_{1} > s_{2} > s_{3}$}
{$s_{1} < s_{2} < s_{3}$}
{$s_{1} = s_{3} > s_{2}$}
{$s_{1} = s_{3} < s_{2}$}
%% answer: D
%% memo:
\memoanswer{
本题原卷及材料中未提供详解，待补充。
}

\item
%% number: 10
%% paperName: 1992年湖南云南海南高考第 10 题 3 分
%% typeId: 1
%% type: 单选题
%% chapter: 直线运动
%% point: 打点计时器公式
%% method: 无
%% score: 3
%% degree: 750
%% duplicateId: 0
%% body:
一同学在用打点计时器做测匀变速直线运动加速度的实验时，纸带上打出的不是圆点，而是如图所示的一些短线，这可能是因为\xzanswer{D}
%% MISSING-IMAGE: 题干图（原 JSON 无 <img>），缺图待补；补图后用 \onepicture 插入
\fourchoices[answer=D]
{打点计时器错接在直流电源上}
{电源电压不稳定}
{电源的频率不稳定}
{打点针压得过紧}
%% answer: D
%% memo:
\memoanswer{
本题原卷及材料中未提供详解，待补充。
}

\item
%% number: 11
%% paperName: 1992年湖南云南海南高考第 11 题 5 分
%% typeId: 2
%% type: 多选题
%% chapter: 曲线运动
%% point: 向心加速度向心力
%% method: 无
%% score: 5
%% degree: 850
%% duplicateId: 0
%% body:
一质点作圆周运动，速度处处不为零。则\xzanswer{ABD}
\fourchoices[answer=ABD]
{任何时刻质点所受的合力一定不为零}
{任何时刻质点的加速度一定不为零}
{质点的速度大小一定不断地改变}
{质点的速度方向一定不断地改变}
%% answer: ABD
%% memo:
\memoanswer{
本题原卷及材料中未提供详解，待补充。
}

\item
%% number: 12
%% paperName: 1992年湖南云南海南高考第 12 题 5 分
%% typeId: 2
%% type: 多选题
%% chapter: 曲线运动
%% point: 人造地球卫星
%% method: 无
%% score: 5
%% degree: 650
%% duplicateId: 0
%% body:
甲乙两颗人造地球卫星，质量相同，它们的轨道都是圆。若甲的运行周期比乙大，则\xzanswer{ACD}
\fourchoices[answer=ACD]
{甲距地面的高度一定比乙大}
{甲的速度一定比乙大}
{甲的加速度一定比乙小}
{甲的动能一定比乙小}
%% answer: ACD
%% memo:
\memoanswer{
本题原卷及材料中未提供详解，待补充。
}

\item
%% number: 13
%% paperName: 1992年湖南云南海南高考第 13 题 5 分
%% typeId: 2
%% type: 多选题
%% chapter: 电场
%% point: 电容
%% method: 无
%% score: 5
%% degree: 650
%% duplicateId: 0
%% body:
一平行板电容器，始终与电池相连。现将一块均匀的电介质板插进电容器，恰好充满两极板间的空间。与未插电介质时相比，\xzanswer{AB}
\fourchoices[answer=AB]
{电容器所带的电量增大}
{电容器的电容增大}
{两极板间各处电场强度减小}
{两极板间的电势差减小}
%% answer: AB
%% memo:
\memoanswer{
本题原卷及材料中未提供详解，待补充。
}

\item
%% number: 14
%% paperName: 1992年湖南云南海南高考第 14 题 5 分
%% typeId: 2
%% type: 多选题
%% chapter: 机械振动
%% point: 振动图象
%% method: 无
%% score: 5
%% degree: 650
%% duplicateId: 0
%% body:
一质点作简谐振动，其位移 $x$ 与时间 $t$ 的关系曲线如图。由图可知\xzanswer{BC}
%% MISSING-IMAGE: 题干图（原 JSON 无 <img>），缺图待补；补图后用 \onepicture 插入
\fourchoices[answer=BC]
{质点振动的频率是 $4\UHz$}
{质点振动的振幅是 $2\Ucm$}
{在 $t = 3$ 秒时，质点的速度为最大}
{在 $t = 4$ 秒时，质点所受的合外力为零}
%% answer: BC
%% memo:
\memoanswer{
本题原卷及材料中未提供详解，待补充。
}

\item
%% number: 15
%% paperName: 1992年湖南云南海南高考第 15 题 5 分
%% typeId: 2
%% type: 多选题
%% chapter: 气体性质
%% point: 气体图象
%% method: 无
%% score: 5
%% degree: 650
%% duplicateId: 0
%% body:
一定质量的理想气体，封闭在带活塞的气缸中，气体从状态 a 出发，经历 ab、bc、cd、da 四个过程回到状态 a，各过程的压强 $p$ 与温度 $T$ 的关系如图所示。其中气体不对外界作功，外界也不对气体作功的过程是\xzanswer{AC}
%% MISSING-IMAGE: 题干图（原 JSON 无 <img>），缺图待补；补图后用 \onepicture 插入
\fourchoices[answer=AC]
{ab 过程}
{bc 过程}
{cd 过程}
{da 过程}
%% answer: AC
%% memo:
\memoanswer{
本题原卷及材料中未提供详解，待补充。
}

\item
%% number: 16
%% paperName: 1992年湖南云南海南高考第 16 题 5 分
%% typeId: 2
%% type: 多选题
%% chapter: 电路
%% point: 电流强度
%% method: 建立模型
%% score: 5
%% degree: 550
%% duplicateId: 0
%% body:
有一横截面积为 $S$ 的铜导线，流经其中的电流强度为 $I$，设每单位体积的导线中有 $n$ 个自由电子，电子的电量为 $q$。此时电子的定向移动速度为 $v$，在 $\Delta t$ 时间内，通过导线横截面的自由电子数目可表示为\xzanswer{AC}
\fourchoices[answer=AC]
{$nvS\Delta t$}
{$nvS\Delta t$}
{$\frac{I\Delta t}{q}$}
{$\frac{I\Delta t}{Sq}$}
%% answer: AC
%% memo:
\memoanswer{
本题原卷及材料中未提供详解，待补充。
}

\item
%% number: 17
%% paperName: 1992年湖南云南海南高考第 17 题 5 分
%% typeId: 1
%% type: 单选题
%% chapter: 电磁感应
%% point: 验证电磁感应实验
%% method: 无
%% score: 5
%% degree: 650
%% duplicateId: 0
%% body:
某学生做观察电磁感应现象的实验，将电流表、线圈 A 和 B、蓄电池、电键用导线连接成如图的实验电路。当他接通、断开电键时，电表的指针都没有偏转，其原因是\xzanswer{A}
%% MISSING-IMAGE: 题干图（原 JSON 无 <img>），缺图待补；补图后用 \onepicture 插入
\fourchoices[answer=A]
{电键位置接错}
{电流表的正负极接反}
{线圈 B 的接头 3、4 接反}
{蓄电池的正、负极接反}
%% answer: A
%% memo:
\memoanswer{
本题原卷及材料中未提供详解，待补充。
}

\item
%% number: 18
%% paperName: 1992年湖南云南海南高考第 18 题 5 分
%% typeId: 2
%% type: 多选题
%% chapter: 机械波
%% point: 波的多解性
%% method: 无
%% score: 5
%% degree: 450
%% duplicateId: 0
%% body:
a、b 是水平绳上的两点，相距 $42\Ucm$。一列正弦横波沿此绳传播，传播方向从 a 到 b。每当 a 点经过平衡位置向上运动时，b 点正好到达上方最大位移处。此波的波长可能是\xzanswer{CD}
\fourchoices[answer=CD]
{$168\Ucm$}
{$84\Ucm$}
{$56\Ucm$}
{$24\Ucm$}
%% answer: CD
%% memo:
\memoanswer{
本题原卷及材料中未提供详解，待补充。
}

\item
%% number: 19
%% paperName: 1992年湖南云南海南高考第 19 题 5 分
%% typeId: 2
%% type: 多选题
%% chapter: 光学
%% point: 全反射
%% method: 无
%% score: 5
%% degree: 650
%% duplicateId: 0
%% body:
图中方框区域内有一个位置可任意摆放的全反射棱镜，其横截面是等腰直角三角形。光线 1、2、3、4 所表示的入射光束经此棱镜后，相应的出射光线是 1$'$、2$'$、3$'$、4$'$。下图 4 种情形中，哪些是可能实现的？\xzanswer{ABC}
%% MISSING-IMAGE: 题干图（原 JSON 无 <img>，四个图片选项缺失），缺图待补；补图后用 \fourchoices[ispicture=true] 插入
%% answer: ABC
%% memo:
\memoanswer{
本题原卷及材料中未提供详解，待补充。
}

\item
%% number: 20
%% paperName: 1992年湖南云南海南高考第 20 题 4 分
%% typeId: 3
%% type: 填空题
%% chapter: 光学
%% point: 透镜
%% method: 无
%% score: 4
%% degree: 850
%% duplicateId: 0
%% body:
做测量凸透镜焦距的实验时，当透镜置于蜡烛与光屏距离的中点时，光屏上恰好得到烛焰的清晰的像，这时测得蜡烛与光屏之间的距离为 $L\Ucm$，则透镜的焦距是 \tkanswer[2cm]{$\frac{L}{4}$}$\Ucm$。
%% answer: 
\jdanswer{
$\frac{L}{4}$
}
%% memo:
\memoanswer{
本题原卷及材料中未提供详解，待补充。
}

\item
%% number: 21
%% paperName: 1992年湖南云南海南高考第 21 题 4 分
%% typeId: 3
%% type: 填空题
%% chapter: 原子和原子核
%% point: 人工核反应
%% method: 无
%% score: 4
%% degree: 850
%% duplicateId: 0
%% body:
有一核反应，其反应式为 $\alpha$ + Be 核 $\to$ n + \ce{^{12}_{6}C}，则 Be 核的质量数是 \tkanswer[1.5cm]{9}，电荷数是 \tkanswer[1.5cm]{4}。
%% answer: 
\jdanswer{
9，4
}
%% memo:
\memoanswer{
本题原卷及材料中未提供详解，待补充。
}

\item
%% number: 22
%% paperName: 1992年湖南云南海南高考第 22 题 4 分
%% typeId: 3
%% type: 填空题
%% chapter: 交流电
%% point: 变压器
%% method: 无
%% score: 4
%% degree: 850
%% duplicateId: 0
%% body:
一理想变压器的原线圈为 3300 匝，接到 $110\UV$ 的电源上，设副线圈输出电压 $5.5\UV$，电流 $20\UmA$。则副线圈的匝数等于 \tkanswer[1.5cm]{165}，原线圈中的电流等于 \tkanswer[1.5cm]{1}$\UmA$。
%% answer: 
\jdanswer{
165，1
}
%% memo:
\memoanswer{
本题原卷及材料中未提供详解，待补充。
}

\item
%% number: 23
%% paperName: 1992年湖南云南海南高考第 23 题 4 分
%% typeId: 3
%% type: 填空题
%% chapter: 光学
%% point: 光子说
%% method: 建立模型
%% score: 4
%% degree: 450
%% duplicateId: 0
%% body:
已知每秒钟从太阳射到地球上垂直于太阳光的每平方米截面上的辐射能为 $1.4\times10^{3}\UJ$，其中可见光部分约占 45\%。假如认为可见光的波长均为 $0.55\Uum$，太阳向各个方向的辐射是均匀的，日地间距离 $R = 1.5\times10^{11}\Um$，普朗克常数 $h = 6.6\times10^{-34}\UJcdots$。由此可估算出太阳每秒钟辐射出的可见光的光子数约为 \tkanswer[2cm]{$4.9\times10^{44}$}个。（只要求二位有效数字）
%% answer: 
\jdanswer{
$4.9\times10^{44}$
}
%% memo:
\memoanswer{
本题原卷及材料中未提供详解，待补充。
}

\item
%% number: 24
%% paperName: 1992年湖南云南海南高考第 24 题 5 分
%% typeId: 4
%% type: 实验题
%% chapter: 电路
%% point: 电表改装
%% method: 无
%% score: 5
%% degree: 750
%% duplicateId: 0
%% body:
一量程为 $100\UuA$ 的电流表，内阻为 $100\UO$，现串联一个 $9900\UO$ 的电阻将它改装成电压表。该电压表的量程是 \tkanswer[1.5cm]{1}$\UV$。用它来测量电压，表盘指针位置如图所示。该电压的大小是 \tkanswer[1.5cm]{0.85}$\UV$。
%% MISSING-IMAGE: 题干图（原 JSON 无 <img>，表盘指针图缺失），缺图待补；补图后用 \onepicture 插入
%% answer: 
\jdanswer{
1，0.85
}
%% memo:
\memoanswer{
本题原卷及材料中未提供详解，待补充。
}

\item
%% number: 25
%% paperName: 1992年湖南云南海南高考第 25 题 5 分
%% typeId: 4
%% type: 实验题
%% chapter: 电路
%% point: 电表改装
%% method: 无
%% score: 5
%% degree: 650
%% duplicateId: 0
%% body:
电流表与电阻箱的某一给定电阻串联，使之改装成一电压表。现要对此电压表进行校准，所用的电路原理如左图所示。右图给出了实验器材的实物图，请按原理图要求连成实验电路。
%% MISSING-IMAGE: 题干图（原 JSON 无 <img>，原理图与实物连线图缺失），缺图待补；补图后用 \onepicture / \drawpicanswer 插入
%% answer: 
\jdanswer{
如图（原卷未提供答案连线图）
}
%% memo:
\memoanswer{
本题原卷及材料中未提供详解，待补充。
}

\item
%% number: 26
%% paperName: 1992年湖南云南海南高考第 26 题 8 分
%% typeId: 6
%% type: 计算题
%% chapter: 运动和力
%% point: 牛顿第二定律
%% method: 无
%% score: 8
%% degree: 650
%% duplicateId: 0
%% body:
将质量为 $m$ 的木块放在倾角为 $\theta$ 的斜面上，木块可沿斜面匀速下滑。现用一沿斜面的力 $F$ 作用于木块，使之沿斜面向上作匀加速运动，如图所示。求木块的加速度。
%% MISSING-IMAGE: 题干图（原 JSON 无 <img>），缺图待补；补图后用 \onepicture[align=right] 插入
%% answer: 
\jdanswer{
$a = \frac{F}{m} - 2g\sin\theta$
}
%% memo:
\memoanswer{
本题原卷及材料中未提供详解，待补充。
}

\item
%% number: 27
%% paperName: 1992年湖南云南海南高考第 27 题 12 分
%% typeId: 6
%% type: 计算题
%% chapter: 电路
%% point: 串并联电路
%% method: 无
%% score: 12
%% degree: 550
%% duplicateId: 0
%% body:
在图示的电路中，电源和电流表的内阻均可不计。当两个电键 K$_{1}$、K$_{2}$ 都断开或都闭合时，电流表 A 的读数是相同的。求电阻 $R$ 的阻值。
%% MISSING-IMAGE: 题干图（原 JSON 无 <img>），缺图待补；补图后用 \onepicture[align=right] 插入
%% answer: 
\jdanswer{
$R = 50\UO$
}
%% memo:
\memoanswer{
本题原卷及材料中未提供详解，待补充。
}

\item
%% number: 28
%% paperName: 1992年湖南云南海南高考第 28 题 14 分
%% typeId: 6
%% type: 计算题
%% chapter: 运动和力
%% point: 动量和能量
%% method: 无
%% score: 14
%% degree: 450
%% duplicateId: 0
%% body:
质量分别为 $m$ 和 $M$ 的两个粒子发生碰撞，碰撞前后两粒子都在同一直线上。在碰撞过程中损失的动能为定值 $E_{0}$。今要求碰撞前两粒子的总动能为最小，求碰撞前两粒子的速度大小和方向。
%% answer: 
\jdanswer{
$v_{0} = \sqrt{\frac{2ME_{0}}{m(M+m)}}$，$V_{0} = -\sqrt{\frac{2mE_{0}}{M(M+m)}}$；碰撞前两粒子速度方向相反。
}
%% memo:
\memoanswer{
本题原卷及材料中未提供详解，待补充。
}

\item
%% number: 29
%% paperName: 1992年湖南云南海南高考第 29 题 15 分
%% typeId: 6
%% type: 计算题
%% chapter: 气体性质
%% point: 气体连接体
%% method: 无
%% score: 15
%% degree: 350
%% duplicateId: 0
%% body:
如右图，A、B 是两个截面积相同的气缸，放在水平地面上，活塞可无摩擦地上、下移动，活塞上固定一细的刚性推杆，顶在一可绕水平固定轴 O 自由旋转的杠杆 MN 上，接触点光滑。活塞（连推杆）、杠杆的质量均可忽略。开始时，A 和 B 中气体压强为 $p_{A} = 1.10\times10^{5}\UPa$ 和 $p_{B} = 1.20\times10^{5}\UPa$。体积均为 $V_{0} = 1.00\UL$，温度均为 $T_{0} = 300\UK$，杠杆处于水平位置。设大气压强始终为 $p_{0} = 1.00\times10^{5}\UPa$，当气缸 B 中气体的温度 $T_{B}$ 变为 $400\UK$、体积 $V_{B} = 1.10\UL$ 时，求气缸 A 中气体的温度。
%% MISSING-IMAGE: 题干图（原 JSON 无 <img>），缺图待补；补图后用 \onepicture[align=right] 插入
%% answer: 
\jdanswer{
$T_{A} = 268\UK$
}
%% memo:
\memoanswer{
本题原卷及材料中未提供详解，待补充。
}

\end{enumerate}

\end{document}
